\section{Discussion and limitations}
\label{sec:disc}

\paragraph{What the evidence supports.}
Source-model results improve the prediction of a new model's accuracy from a few items when the new model resembles the sources (random split): by about 1.2--1.6 MAE points at $n{=}50$ over a random subset on the macro-average, with the largest gains at the smallest budgets and for ranking. They do not by themselves protect against capability shift: when the targets are stronger than all sources, methods that average the predictions of a fitted model of the sources are systematically biased downward and lose to random sampling, in both MAE and rank correlation. A design-based residual correction (the difference estimator of survey sampling) added to a fitted model removes that bias: the ablation shows that it, and not stratification, is what protects the model-based predictor under extrapolation, that random items suffice for it, and that stratification adds a further gain mostly when the estimator is the plain subset mean. Corrected estimators retain a gain over a random subset of about $0.7$--$1.6$ points at $n{=}50$ in both regimes (not significantly worse than random in any benchmark, split and budget we tested with the target-model bootstrap), but the correction costs accuracy when the sources are informative (random split: $3.06$ for the plain stratified model against $3.40$ with the correction), so an uncorrected or shrunk model is better there. The benefit of stratification on top of the correction is $0.12$ points at $n{=}50$ at most. No method ranks models better than a random subset under extrapolation, as far as we can detect. We do not claim that any method is optimal, only the ordering shown on this data and protocol.

\paragraph{Limitations.}
(i) \emph{Re-implementations.} The anchor-point and IRT pipelines are our versions (a $k$-means on an SVD embedding, MAP 2PL-IRT without gp-IRT mixing, fixed dimensions), not the original code, and the original methods use confidence information or large model pools that binary matrices of 79--448 models lack; their numbers here need not match the published ones, and the original packages were not run. Active Evaluation Acquisition, the metabench GAM predictor and the item-feature assessor are missing. (ii) \emph{Hyperparameters.} IRT and low-rank dimensions were fixed, with one informal check of three IRT dimensions on three benchmarks during debugging (Appendix~\ref{app:honest}), so the model-based baselines might improve with tuning on source models (we did not test whether this removes their extrapolation bias). (iii) \emph{Targets and benchmarks.} There are 20--107 target models per benchmark, they come from two leaderboards of a particular era, and the benchmarks share models, so the 13 benchmarks are not independent replicates. The extrapolation split is defined by accuracy of the model on the same benchmark; other shifts (a different family, newer models, other prompts) were not tested. (iv) \emph{Budgets.} $n\le200$ and $n\le N/4$; statements about the number of items needed are limited by this cap. (v) \emph{Independence of the design of P3.} Its candidates were chosen after a handful of pilot runs and after the implementer had read the two-sentence summary in the data repository's README (Appendix~\ref{app:honest}); the final choice rule was fixed in advance and uses source models only. (vi) \emph{Ranking.} Rank correlations over a handful of targets are noisy and depend on the score spread; the pairwise analysis shows that models closer than about 2 points are barely ordered above chance by any method. (vii) The further analyses (Section~\ref{sec:exp}) use 4--7 benchmarks and 20 repetitions, and the organisation hold-out is not a capability or family shift of the kind in the plan.

\paragraph{Broader impact.}
Cheaper evaluation lowers cost, but a few-item estimate can hide failures elsewhere in the benchmark, rank close models arbitrarily and invite gaming of the selected items; our results warn against trusting source-model-based predictions for models stronger than the sources.

\paragraph{Use of LLM agents.}
The study, including code, experiments, figures and text, was carried out by LLM-based agents under human-specified objectives; the estimators do not use an LLM. Every number is generated by the delivered scripts, and the cited works were consulted at the level of their abstracts, not read in full.
